\section{Related Work}
\label{sec:related}

\paragraph{FAIR and its descendants.}
The FAIR Guiding Principles are the
substrate~\cite{wilkinson2016fair}. FAIR4RS~\cite{chuehong2022fair4rs}
re-reads them for research software; FAIR4ML~\cite{ravi2024fair4ml}
re-reads them for machine-learning models.
F(AI)\textsuperscript{2}R is the third such re-reading and the first
that targets the artefact classes generated downstream of the writing
process itself. The relationship is additive, not revisionist: a paper
whose data is FAIR, whose code is FAIR4RS, and whose models are
FAIR4ML is, on the writing layer,
\emph{FAIR-uncovered until it is also F(AI)\textsuperscript{2}R}.

\paragraph{Datasheets and model cards.}
Datasheets for datasets~\cite{gebru2021datasheets} and model cards for
model reporting~\cite{mitchell2019modelcards} are the canonical
attempts to attach human-readable, structured documentation to ML
artefacts. Both are ancestors of the per-prompt and per-claim
documentation we propose. They differ in scope --- datasheets and
model cards document an artefact at a single level; F(AI)\textsuperscript{2}R
documents an artefact \emph{and} the process that produced it,
chained.

\paragraph{Provenance and reproducibility.}
PROV-O~\cite{w3c2013provo} is the formal substrate. The recent
ML-reproducibility literature has produced a number of tools
(\textsc{papers-with-code} reproducibility checklists; the NeurIPS
reproducibility checklist; \textsc{Code Ocean}-style executable
papers). F(AI)\textsuperscript{2}R differs in that the writing layer
is itself the artefact under reproducibility audit, not merely the
code or the dataset.

\paragraph{Systematic reviews and certainty of evidence.}
PRISMA 2020~\cite{page2021prisma} and GRADE~\cite{guyatt2008grade}
supply the retrieval-depth and invocation-strength axes the
verification ladder crosswalks. Neither was designed for LLM-assisted
writing; both have decades of editorial practice behind them. The
ladder is, in this sense, an attempt to import established practice
from systematic-review methodology into a setting that has lacked it.

\paragraph{AI-disclosure norms and authorship.}
The \textsc{ICMJE}~\cite{icmje2023}, ACL Rolling Review, NeurIPS, and
the major scientific publishers have converged on a position that
disclosure is mandatory and AI authorship is unacceptable. The
literature on what \emph{adequate} disclosure looks like is thinner.
Bender and colleagues' critiques~\cite{bender2021stochastic} provide
much of the vocabulary the disclosure community has since used; the
sixth practice (\S\ref{sec:practice-disclosure}) is our attempt to
operationalise it in a single submission.

\paragraph{The base-rate literature on fabrication.}
Reported base rates for fabricated citations in LLM
output~\cite{walters2023fabrication,magesh2024legal} are the load-bearing
empirical input to the verification ladder. Without a number, the
ladder is overhead; with a number, the ladder is the cheapest known
mitigation that scales.

\paragraph{Recursive-meta-process examples.}
A handful of recent papers describe their own LLM-assisted production
process. We are not aware of a prior recursive case study that
combines transcript preservation, a verification-status FSM, a PROV-O
graph, and the full eight-practice integration we describe here. The
source repository \texttt{noheton/Obscurity-Is-Dead} is the immediate
methodological ancestor of this paper; this paper is its abstraction.
